\documentclass[11pt,a4paper]{article}
\usepackage[margin=2.2cm]{geometry}
\usepackage[T1]{fontenc}
\usepackage{lmodern}
\usepackage{amsmath,amssymb}
\usepackage{booktabs}
\usepackage{xcolor}
\usepackage{enumitem}
\usepackage{titlesec}
\usepackage[colorlinks=true,linkcolor=black,urlcolor=blue]{hyperref}

\definecolor{revblue}{rgb}{0,0,0.75}
\definecolor{revred}{rgb}{0.75,0,0}
\definecolor{boxgrey}{gray}{0.94}

\titleformat{\section}{\large\bfseries}{\thesection.}{0.6em}{}
\titlespacing{\section}{0pt}{14pt}{6pt}
\titleformat{\subsection}{\bfseries}{}{0em}{}
\titlespacing{\subsection}{0pt}{10pt}{3pt}

\newsavebox{\calloutbox}
\newenvironment{callout}[1]
  {\par\medskip\noindent\begin{lrbox}{\calloutbox}%
   \begin{minipage}{\dimexpr\linewidth-20pt\relax}%
   \textbf{#1}\par\smallskip}
  {\end{minipage}\end{lrbox}%
   \setlength{\fboxsep}{10pt}%
   \colorbox{boxgrey}{\usebox{\calloutbox}}\par\medskip}

\setlist[itemize]{leftmargin=1.3em,itemsep=2pt,topsep=3pt}
\newcommand{\blue}[1]{\textcolor{revblue}{#1}}

\begin{document}

\begin{center}
{\LARGE\bfseries Meeting Brief}\\[3pt]
{\large The blue revision of \emph{Two Portable Components for Meeting Hard
Prediction Quotas}}\\[6pt]
{\small Prepared 9 August 2026 \quad$\cdot$\quad read this one document; everything else is backup}
\end{center}

\vspace{4pt}
\hrule
\vspace{10pt}

\section{What you are presenting, in four sentences}

The manuscript is unchanged in substance: same claims, same headline numbers, same
structure. What is new is a body of follow-up work (Track B, roughly 600 runs) that
answers the open questions the paper itself flagged, and five insertions that fold those
answers in. \blue{Every insertion is blue.} \textcolor{revred}{Every sentence the new work
made obsolete is struck in red with a blue replacement beside it.} Your professor can read
the paper start to finish and accept the revision without editing a word of his own text.

\begin{callout}{The single most important fact to state up front}
\texttt{paper/main.tex} --- the version he wrote --- is byte-identical to what he sent.
The revision lives in a separate file. This was verified mechanically, not by eye: a
checker strips every blue and red character and compares the remainder to his original at
character level. It reports that not one character of his prose was added, removed, or
altered. The only structural change is one figure relocated from \S1 to Appendix E.
\end{callout}

\section{The five blue insertions}

\begin{enumerate}[leftmargin=1.5em]
\item \textbf{\S4, one clause} --- notes that the native-resolution question is no longer
entirely deferred. Sets up insertion 4.

\item \textbf{\S5.1, one paragraph + Table 2} --- three robustness checks:
a stronger dual baseline (ALM), a fourth backbone (MobileNetV2), and interval estimates
(bootstrap CIs, BH-FDR). \emph{From B3, B8, B6, B7.}

\item \textbf{\S5.3, one paragraph + Table 3} --- the imbalanced-learning baselines
(focal, class-balanced, logit adjustment), each trained with its own loss and then
clipped. \emph{From B1.}

\item \textbf{\S7, one paragraph + Table 4} --- a native-resolution check on HAM10000, one
dataset run to full depth: every method, every backbone, two caps, four seeds, 96 runs.
\emph{From B2, reworked.}

\item \textbf{Ablation appendix, one paragraph} --- the component ablation repeated in the
tie region. \emph{From B4.}
\end{enumerate}

\section{Track B, one line each}

\begin{itemize}
\item \textbf{B1 --- imbalanced baselines} (312 runs). Focal is the only competitive
member. The pattern is \emph{dataset}-driven: on DermMNIST focal matches or slightly beats
TraLO's macro-F1; elsewhere TraLO leads or ties. None of the three satisfies the cap
without the post-hoc LP.
\item \textbf{B2 --- native resolution} (96 runs, native HAM10000). The macro-F1 margin
over clipping reproduces: \textbf{+0.013, positive in all 6 backbone$\times$cap cells}.
cc-F1 ties both duals. Satisfaction: TraLO 24/24, Hounie 24/24, Fioretto 0.96, LP-LG 0/24.
\item \textbf{B3 --- ALM} (24 runs). \emph{The strongest new result.} ALM is genuinely a
stronger dual than Fioretto (beats it by +0.010), and TraLO still tops it in
\textbf{all six} tight-cap cells by +0.028 cc-F1.
\item \textbf{B4 --- tie-region ablation} (20 runs). The reset and hinge carry the
\emph{tight-cap} advantage specifically, not a general one. Consistent with the mechanism:
they bind only where the cap binds.
\item \textbf{B5 --- win-bar sensitivity} (no new runs). The dual win at L30/L40 is stable
at every threshold (4/0/0 even at the strict $\tau{=}0.010$). The clip comparison never
produces a stable win at any threshold.
\item \textbf{B6 --- bootstrap CIs} (no new runs). 20\,000 resamples on the tight-cap
cells.
\item \textbf{B7 --- BH-FDR} (no new runs). 3 of 8 cells survive against Fioretto; 0 of 8
against the clipper.
\item \textbf{B8 --- MobileNetV2} (32 runs). The advantage reproduces on a fourth
backbone.
\end{itemize}

\begin{callout}{The through-line, if he asks what Track B established}
The advantage is \emph{specific}, and the follow-ups sharpen where. Against the
\textbf{dual-ascent family} at tight caps it is real, survives a stronger dual (ALM),
reproduces on a fourth backbone, and holds at every win threshold. Against
\textbf{post-hoc clipping} the durable, categorical separation is \emph{native cap
satisfaction} (1.0 vs 0.0 everywhere) plus a macro-F1 margin that now also reproduces at
native resolution. The honest boundary: on the imbalance-friendly dataset a focal clipper
closes the macro-F1 gap. That boundary is now stated in the paper rather than left open.
\end{callout}

\section{The three imbalanced-learning baselines (B1)}

First, the framing point, because it is the one most easily got wrong: \textbf{these are not
clipping algorithms.} They are \emph{training losses} --- three ways of making a classifier
care more about a rare class while it learns. Clipping is a separate step applied
afterwards. That is exactly why they belong in the paper: none of the three controls
\emph{how many} items get the constrained label, so each still needs the post-hoc LP to
meet the cap. Their presence tests whether a smarter \emph{training} recipe plus clipping
can match constraint-time training.

All three run at their original papers' settings, verified against the source texts.

\subsection{Focal loss \quad\small\texttt{lin2017focal} --- Lin et al., ICCV 2017}
\[
\mathcal{L}_{\text{focal}} \;=\; -\,\alpha\,(1-p_y)^{\gamma}\,\log p_y ,
\qquad \alpha = 0.25,\;\; \gamma = 2
\]
where $p_y$ is the predicted probability of the \emph{true} class.

\begin{itemize}
\item $\gamma$ is the \textbf{focusing exponent}. It decides how fast already-correct
examples stop mattering. At $p_y{=}0.9$, the factor $(1-0.9)^2 = 0.01$ --- that example
contributes $100\times$ less than under plain CE. At $p_y{=}0.5$ it is only $4\times$ less.
\item $\alpha$ is a fixed scale on the whole term.
\end{itemize}

\textbf{Why it suits imbalanced data:} the majority class is learned early and becomes
``easy''. Under plain CE those easy examples still dominate the gradient by sheer count.
Focal drains their contribution so learning stays pointed at what is still hard --- which
is disproportionately the rare class. Note it reweights by \emph{difficulty}, which
correlates with rarity but is not the same thing.

\subsection{Class-balanced reweighting \quad\small\texttt{cui2019class} --- Cui et al., CVPR 2019}
\[
E_{n_c} \;=\; \frac{1-\beta^{\,n_c}}{1-\beta}
\qquad\Longrightarrow\qquad
w_c \;\propto\; \frac{1}{E_{n_c}} \;=\; \frac{1-\beta}{1-\beta^{\,n_c}},
\qquad \beta = 0.9999
\]
with $n_c$ the number of \emph{training} samples of class $c$; weights are normalised to
mean 1 so the loss scale still matches CE.

\begin{itemize}
\item $\beta$ sets a \textbf{saturation scale} of $1/(1-\beta) = 10{,}000$ samples. Below
that, a class is weighted close to inverse-frequency; far above it, extra samples stop
adding weight.
\end{itemize}

\textbf{Why it suits imbalanced data:} the naive fix, weighting by $1/n_c$, over-corrects.
Cui et al.'s argument is that samples overlap --- the 40{,}000th image of a class teaches
almost nothing new --- so the useful quantity is an \emph{effective} number, not a raw
count. \textbf{Why it is active on our data:} OctMNIST's training classes run from roughly
$7.7$k to $46$k, straddling the $10$k saturation scale, so some classes sit below the
saturation point and others well above --- the weights genuinely differ.

\subsection{Logit adjustment \quad\small\texttt{menon2020long} --- Menon et al., ICLR 2021}
\[
\mathcal{L}_{\text{LA}} \;=\; \mathrm{CE}\!\left(z + \tau\log\pi,\; y\right),
\qquad \tau = 1
\]
where $z$ are the logits and $\pi$ the training class prior.

\begin{itemize}
\item $\tau$ scales the prior correction. $\tau{=}1$ is not a tuned knob: it is the value
that makes the classifier Bayes-optimal for \emph{balanced} error. Anything else would need
justifying.
\end{itemize}

\textbf{Why it suits imbalanced data:} a model trained on a skewed set learns the training
prior along with the signal. Adding $\tau\log\pi$ to the logits \emph{during} training
handicaps the frequent classes, so the network must earn its confidence on them instead of
coasting on base rates. It shifts the decision boundary rather than reweighting anything.

\begin{callout}{Why these three, and not three flavours of one idea}
They attack three genuinely different mechanisms: focal reweights per \emph{example} by
difficulty; class-balanced reweights per \emph{class} by frequency; logit adjustment moves
the \emph{decision boundary} by the prior. That makes them a fair panel rather than a
stacked one. \textbf{If asked whether they were handicapped:} they run at their authors'
own settings, and focal \emph{beats} TraLO on DermMNIST --- these are not baselines set up
to fail.
\end{callout}

\section{Reading the statistics (B5--B7)}

Three different ways of interrogating the same question: \emph{is the OctMNIST tight-cap
cc-F1 result real, or an artifact of a small sample?} Each guards against a different
failure. All three are re-analyses --- no new runs.

\subsection{The underlying problem: $n=4$ seeds}
Four seeds per cell is very few, and it creates two separate worries. \emph{Is this one
number reliable?} (B6.) \emph{I ran many comparisons --- how many won by luck?} (B7.)
Fixing one does not fix the other. B5 adds a third: \emph{does my conclusion depend on
where I drew the line?}

\subsection{B5 --- win-bar sensitivity}
The paper calls a cell a ``win'' when the gap clears $+0.005$. That threshold is a
judgement call, so B5 re-scores everything at $\tau \in \{0.003, 0.005, 0.010\}$.
\textbf{Result:} against Fioretto the tight-cap win is 4/0/0 at \emph{every} threshold,
including the strict $0.010$. Against the clipper it is never stable at any threshold.
\textbf{What it protects against:} threshold-shopping. A conclusion that flips when you
nudge the bar is a property of the bar, not of the data.

\subsection{B6 --- bootstrap confidence intervals}
With $n{=}4$ you cannot lean on a $t$-test's normality assumption. Instead, resample the
paired differences \emph{with replacement} 20{,}000 times, recompute the mean each time,
and take the middle 95\% of those means. That range is the interval --- an answer to ``how
much would this number move on a slightly different sample?'' without assuming a
distribution. The resampling unit is \emph{cells}, not seeds, which respects the
atomic-unit rule.

\begin{center}\small
\begin{tabular}{llrcl}
\toprule
comparison & cap & mean & 95\% CI & reading \\
\midrule
TraLO $-$ Fioretto & L30 & $+0.046$ & $[+0.024, +0.074]$ & entirely above 0 \\
TraLO $-$ Fioretto & L40 & $+0.033$ & $[+0.021, +0.044]$ & entirely above 0 \\
TraLO $-$ ALM      & L30 & $+0.035$ & $[+0.010, +0.048]$ & entirely above 0 \\
TraLO $-$ clip     & L30 & $-0.003$ & $[-0.013, +0.008]$ & \textbf{straddles 0} \\
TraLO $-$ clip     & L40 & $+0.002$ & $[-0.010, +0.012]$ & \textbf{straddles 0} \\
\bottomrule
\end{tabular}
\end{center}

\noindent The specific cell reviewers flagged as weak --- ViT-B/16 at L30, about $1.8\sigma$
--- comes out at $+0.086$ against Fioretto and survives. Note also the \emph{shape}: at
L10/L15/L20 even the dual gap ties, because the tightest caps ceiling-crush every method.
The win is a mid-cap phenomenon.

\subsection{B7 --- Benjamini--Hochberg FDR}
Run 20 tests at $p<0.05$ with nothing real going on and you should \emph{expect} about one
``significant'' result. BH controls the \textbf{false discovery rate}: among the results you
declare significant, at most $q{=}5\%$ are expected to be false. Mechanically, the bar gets
stricter the more tests you ran. \textbf{Result:} 3 of 8 cells survive against Fioretto;
\textbf{0 of 8} against the clipper.

\subsection{Why ``3 of 8'' is not the disappointment it looks like}
With 4 seeds, a within-cell paired Wilcoxon test \textbf{cannot return $p<0.125$} --- there
are not enough orderings. That is a floor imposed by the sample size, not by the effect.
Most cells therefore cannot reach individual significance however large the true effect is.
So the meaningful quantity is not 3/8 in isolation; it is \textbf{3 versus 0}.

\begin{callout}{What ``0 of 8 on clipping'' does and does not mean --- know this exactly}
It means: \emph{on cc-F1, at OctMNIST tight caps, TraLO does not beat the post-hoc clipper.}
The paper reports that comparison as a \textbf{tie}, and B5, B6 and B7 are three independent
reasons why.

\smallskip
It does \textbf{not} mean TraLO fails to beat clipping in general. The deployment claim ---
overall quality, macro-F1, $+1.6$ to $+5.3$\,pp, positive in 24 of 27 cells --- is a
\emph{different metric} from a \emph{different analysis}, and it stands. So does native cap
satisfaction, $1.0$ vs $0.0$, which is categorical.

\smallskip
Do not let these two get conflated in the room. B5--B7 are about \emph{cc-F1 at tight
OctMNIST caps} only.
\end{callout}

\subsection{Which of B5--B7 actually reached the manuscript}
Worth knowing precisely, because two are in and one is not.

\begin{center}\small
\begin{tabular}{llp{6.4cm}}
\toprule
& in the paper? & where \\
\midrule
\textbf{B6} bootstrap CIs & \textbf{yes}, in blue & \S5.1 check \emph{(vi)},
\texttt{main.tex:746} --- the $+0.046$ / $+0.033$ gaps with their CIs \\
\textbf{B7} BH-FDR        & \textbf{yes}, in blue & same sentence,
\texttt{main.tex:749} --- ``three of the eight \ldots{} and none against the post-hoc
clipper'' \\
\textbf{B5} win-bar sensitivity & \textbf{no} & --- \\
\bottomrule
\end{tabular}
\end{center}

\noindent B6 and B7 sit in a single long sentence, which is why they can read as one
result rather than two. The other Benjamini--Hochberg mentions (\texttt{main.tex:637--645})
are \emph{pre-existing black text}, not B7.

\begin{callout}{If he asks why B5 was left out}
Not because it was negative --- it is one of the more supportive results in the campaign:
the tight-cap win is 4/0/0 against Fioretto-LDF at \emph{every} threshold tested
($\tau\in\{0.003,0.005,0.010\}$), and the clip comparison is never stable at any of them.
The likely reason is \textbf{redundancy}: the paper already argues threshold-robustness at
\texttt{main.tex:802}, by a different route --- that the cell-mean gaps range $+0.016$ to
$+0.081$, so no plausible threshold changes the classification. That sentence is the
professor's original text.

\smallskip
\textbf{The case for adding it:} B5 upgrades that argument from ``the gaps are large, so
the threshold cannot matter'' to ``we re-scored at three thresholds and the classification
did not change'' --- a directly tested claim rather than an inferred one. One clause, in
blue, and it strengthens an existing sentence rather than introducing a new result.
\end{callout}

\subsection{Where the paper already says this --- exact lines to point at}
The paper uses \textbf{two metrics with two different comparators}, deliberately:
\textbf{cc-F1 is measured against the trained duals; macro-F1 against the clipper.} If he
asks where you concede the clip tie:

\begin{itemize}
\item \texttt{main.tex:651} --- the section is titled ``Constrained-class quality
\emph{vs.\ the trained baselines}''.
\item \texttt{main.tex:659} --- ``the package leads \emph{the best constraint-trained
baseline} on every backbone, by a mean cc-F1 gap of $+0.038$''.
\item \texttt{main.tex:708} --- stronger than a tie, an explicit concession: ``Where a
\emph{post-hoc clipper posts higher raw cc-F1} (three of the six band cells, by up to
$+0.018$), the edge is bought by editing \ldots{} TraLO leads that same cell $0.702$ to
$0.666$ on macro-F1''.
\item \texttt{main.tex:677} --- Figure~1 footnote: ``faded post-hoc clippers can post
higher raw cc-F1 by editing''.
\item \texttt{main.tex:1764} --- appendix column headers: ``$\Delta$cc-F1 vs.\
\emph{trained}'' and ``$\Delta$macro-F1 vs.\ \emph{clipper}''.
\end{itemize}

\noindent So on cc-F1 at tight caps the clipper wins 3 of 6 cells outright. The paper's
answer is not to deny it but to show \emph{how} it is bought --- clipping rewrites roughly
$11\%$ of the batch --- and that TraLO wins the same cell on macro-F1.

\begin{callout}{One wording fix worth raising yourself}
The abstract states this result \emph{without naming the comparator}: ``one regime shows a
consistent advantage across all backbones tested: OctMNIST at tight caps''. Read cold that
scans as an advantage over \emph{everything}; the scoping only arrives in \S5.1. Adding
three words --- ``advantage \emph{over the trained baselines}'' --- removes the ambiguity.
Better to propose this than to have a reviewer read it the loose way.
\end{callout}

\noindent\textbf{Where this helps you.} A null result you found yourself and reported is a
much stronger position than one a reviewer finds for you. You tested the clip comparison
three ways, it did not hold, and the paper says so --- which is also why the headline is
anchored on the dual-ascent comparison, where all three methods agree the effect is real.
It is consistent with the July quota-fill finding: the apparent cc-F1 edge over clipping was
a budget-fill artifact, and the durable separation from clipping is native satisfaction, not
a quality margin on that metric.

\begin{callout}{Warning: two different B6/B7 analyses exist in the repo}
\texttt{results/stats\_trackb/b6/} and \texttt{b7/} are an \emph{older, broader} sweep
covering EfficientNetB0, ResNet18, aider and eurosat --- backbones and datasets not in the
current grid --- reporting 259/270 wins surviving FDR. Those are \textbf{not} the paper's
numbers. The B6/B7 above are the Track-B ones, scoped to the OctMNIST tight-cap family.
\end{callout}

\section{Numbers to know cold}

\begin{center}\small
\begin{tabular}{p{6.6cm}l}
\toprule
Quantity & Value \\
\midrule
Main grid & 1944 runs, 3 datasets $\times$ 3 backbones $\times$ 9 caps $\times$ 4 seeds \\
Track B added & $\approx$600 runs across B1--B8 \\
Headline macro-F1 gain over clipping & $+1.6$ to $+5.3$ pp (dataset means) \\
OctMNIST tight-cap cc-F1 lead vs duals & $+0.038$ over 6 cells, sign test $p{=}0.031$ \\
TraLO $-$ ALM, tight caps & $+0.028$ cc-F1, 6/6 cells \\
ALM $-$ Fioretto & $+0.010$ cc-F1 (so ALM is the tougher baseline) \\
Native HAM10000, macro-F1 vs LP-LG & $+0.013$, 6/6 cells positive \\
Native satisfaction (TraLO / LP-LG) & 24/24 \; vs \; 0/24 \\
Atomic unit of analysis & (dataset, backbone, cap) over 4 seeds --- never pooled \\
\bottomrule
\end{tabular}
\end{center}

\section{What changed last night, and why}

An audit of the full 32 pages extracted 875 checkable assertions and found five real
problems. Two were acted on.

\subsection{Four stale disclaimers}
Four passages written \emph{before} Track B existed now contradicted the blue work --- the
paper said it had not done things it now does. Each is struck in red with a blue
replacement:

\begin{itemize}
\item abstract, intro, \S7: ``native-resolution generalization is left to future work''
\item \S2: ``we do not run a dedicated augmented-Lagrangian baseline'' (ALM is now Table 2)
\item \S2 and \S7: the imbalanced comparison described as not run (it is now \S5.3)
\end{itemize}

\noindent These are all cases where the new work \emph{strengthens} the paper and the old
text had not caught up.

\subsection{Table 3 lost a column}
Table 3 had a \texttt{vs.\ LP-LG} column taken from the B1 campaign's own clip arm. Table~1
reports the same quantity from the main corpus, and the two disagreed in all nine cells ---
including a sign flip, and $+0.001$ where Table~1 says $+0.081$.

\begin{callout}{If he asks about this, the answer is short}
Two different run sets were being asked the same question. The column was deleted, not
recomputed --- substituting corpus numbers into a B1 table would mix run sets inside one
row, which is worse. TraLO-versus-clipping belongs in Table~1 and Appendix~D, over the
full grid. Note the direction: \emph{the deleted column understated the result.} Table~1's
numbers are the larger ones.
\end{callout}

\subsection{Table 4's bolding was fighting its own caption}
The rule was ``best in row'', so Fioretto took bold in 4 of 6 macro-F1 rows --- on margins
of $0.001$--$0.002$, i.e.\ \emph{inside} the $\pm0.005$ band the paper calls a tie. The
table crowned a winner where the caption said there was none. Bold now marks everything
within $0.005$ of the row best, which leaves \textbf{LP-LG as the only unbolded method in
all six macro-F1 rows} --- the claim made visible. The caption also said ``every cell is a
tie'' against the duals; one cell (MobileNetV3, $L40$) leads Hounie-RCL by $+0.008$, a
TraLO win that was being under-reported.

\subsection{The ablation table (Table 5) --- two corrections}
An independent reviewer re-derived four findings that were already sitting unapplied in
\texttt{docs/archive/AUDIT\_2026-07-31.md}. Two were table-level defects and are now fixed
in red/blue:

\begin{itemize}
\item \textbf{The ``Undershoot hinge'' row never isolated the hinge.}
\texttt{src/methodologies/tralo\_bounded/train.py} contains \emph{zero} occurrences of the
optimizer reset, so TraLO $-$ TraLO-bounded removes \emph{both} components. The table's own
numbers give it away: reset alone costs $+0.079$, reset$+$hinge costs $+0.036$ --- removing
more cannot hurt less. Row relabelled \textbf{``Reset $+$ hinge (TraLO-bounded)''}, and the
caption now states that the hinge's individual contribution is not identified \emph{on
these cells}, and records that a correctly isolated hinge arm did run on DermMNIST and
TissueMNIST --- just not on OctMNIST (see Sec.~8).
\item \textbf{A fifth ablation arm was missing.}
\texttt{ablation\_complete.csv} has five arms at $n{=}24$; the table showed four. The
omitted one is \texttt{+KL} (mean $-0.0103$, $p{=}0.0398$) --- significant, and pointing
against the main text's treatment of KL. Added as an explicitly labelled
\emph{addition, not removal} row; the caption no longer claims to be ``Complete''.
\end{itemize}

\noindent Also fixed: Table 9's threshold language (it called $+0.002$ ``a positive
margin'' while calling $-0.003$ elsewhere ``no systematic advantage''), and the
``eight tight-cap cells'' phrasing, which now says \emph{which} eight.

\subsection{Bibliography}
51 citations, 55 entries, zero undefined. Fixed: the Cooper library reference was dated
2024 but was posted April 2025; missing diacritics in Pathak et al.; and ALM had been
introduced with no citation at all (now Bertsekas 1982). The two \emph{added} entries also
render blue in the reference list --- the \texttt{.bib} carried no colour, so they were
appearing black even though their in-text citations were blue. The two metadata
\emph{corrections} above are not marked, since they modify pre-existing entries.

\section{Questions he is likely to ask}

\subsection{``Did you touch my text?''}
No, and it was checked mechanically rather than by eye. Struck sentences are still
character-for-character his --- only a wrapper was added around them, which the checker
confirms.

\subsection{``Why is the figure caption entirely blue?''}
The dataset figure moved from \S1 to Appendix E. An image cannot be colour-marked, so the
relocation is signalled by making its caption blue. His wording is preserved almost
verbatim; one clause was added. \emph{This is a judgement call --- you may prefer to blue
only the new clause and leave his caption text black.}

\subsection{``Does focal beating you on Derm undercut the paper?''}
No, and it is stated in the paper now rather than buried. It narrows the claim from
``constraint-trained beats clipping'' to ``beats \emph{vanilla} clipping'' on that one
dataset. The count-satisfaction separation is untouched: none of the three imbalanced
recipes meets the cap without the post-hoc LP.

\subsection{``Is the cc-F1 win real?''}
Against the duals, yes --- and it is the claim with the most support behind it: stable at
every win threshold (B5), surviving a stronger dual (B3), reproducing on a fourth backbone
(B8), FDR-significant (B7). Against the clipper it is \emph{not} claimed as a win, and B5/B7
say so explicitly (0 of 8 cells survive FDR). The paper reports the clip comparison as a
tie.

\section{Open items --- know these before he raises them}

\begin{itemize}
\item \textbf{The ``1.6--5.3 pp'' range pools across backbones and caps.} Each endpoint is
a dataset mean over $3\times9\times4$ runs. Appendix D's own stated rule is that the atomic
cell is (dataset, backbone, cap) over seeds \emph{only}. At that granularity the same
quantity ranges from $-0.018$ to $+0.091$. The headline is defensible as a summary but it
is not scoped the way the appendix says to scope things. \emph{Unresolved --- your call.}
\item \textbf{Table 1's caption says ``the full symmetric grid''} but shows 3 of 9 caps
(27 of 81 cells), and the L40 number quoted beside it is not in the table.
\item \textbf{Only one of the ``two portable components'' is isolated \emph{on the cells
that carry the claim}.} Following from the Table 5 correction above: on the OctMNIST
tight-cap cells the reset is isolated singly and the hinge is not.

\smallskip
A correctly isolated hinge arm --- \texttt{hybrid\_mode=bounded\_only} with
\texttt{reset\_optimizer\_at\_sat=True}, so the reset is retained --- \textbf{does exist
and did complete}: $7$ runs on DermMNIST and $5$ on TissueMNIST (verified on dsisco01,
2026-08-13). There are \textbf{zero} such runs on OctMNIST, which is the only dataset
Table~5 reports, so it cannot be substituted in.

\smallskip
On the two datasets where it ran ($L30$, MobileNetV3, $3$ seeds), the hinge and the reset
are \emph{each} individually neutral, $\lvert\Delta$~cc-F1$\rvert\le0.007$. That is not a
problem for the paper --- it corroborates the caption's existing caveat that component
importances are estimated inside the win region and need not transfer outside it. The
components bind where the cap binds.

\smallskip
Since ``two portable components'' is billed as the paper's central mechanistic finding,
this is the likeliest place for a reviewer to press. \emph{Cost to settle it: 24 runs
(2 caps $\times$ 3 backbones $\times$ 4 seeds) on OctMNIST. The config already exists, so
it is a dispatch, not new code.}
\item \textbf{OctMNIST has no native-resolution counterpart}, so whether the tight-cap cc-F1
win survives at full resolution is untested. Stated plainly in the paper.
\item Roughly 70 further minor findings from the audit were never adversarially verified ---
the run hit its limit. Treat them as unchecked leads, not defects.
\end{itemize}

\begin{callout}{The one you should expect him to raise: the OctMNIST headline}
Two independent reviews --- the July audit and a fresh one --- have questioned whether the
$+0.038$ cc-F1 lead at OctMNIST tight caps is a \emph{quality} advantage or a
\emph{budget-utilisation} one. The argument runs: precision is flat across the three
trained methods, so the cc-F1 ordering tracks how much of the cap each method spends. The
paper already describes this mechanism in its own words (\S5.1: the duals ``overshoot it
downward and finish $11$--$36\%$ under the budget at $L30$ \ldots{} so recall falls with
the unspent budget while precision holds'') and frames it as TraLO recovering the leftover
headroom.

\smallskip
A reframing was drafted and \textbf{withdrawn at the author's direction}; the abstract and
\S5.1 read exactly as the professor wrote them. Nothing in the paper currently states the
budget-utilisation reading.

\smallskip
\textbf{If he raises it}, the facts to have ready: the decisive control --- refill each
dual's unspent budget with its next-most-confident predictions and re-score --- has not
been run. And the claim is scoped to the trained duals throughout (\S5.1's title and first
sentence), not to clipping, which B5--B7 already report as a tie. The supporting dossier
is \texttt{docs/archive/AUDIT\_2026-07-31.md}.
\end{callout}

\section{Files, if he wants them}

\begin{center}\small
\begin{tabular}{ll}
\toprule
\texttt{paper/main.tex} & his original, untouched \\
\texttt{docs/main.tex} & the blue revision (source) \\
\texttt{paper/main\_rev.pdf} & the blue revision, 32pp --- \textbf{the one to read} \\
\texttt{paper/references.bib} & 55 entries, 51 cited \\
\texttt{docs/track\_b/TRACK\_B\_RESULTS\_HANDOFF.md} & full B1--B8 record \\
\bottomrule
\end{tabular}
\end{center}

\vspace{6pt}
\noindent\small Committed as \texttt{8648b15e} on branch \texttt{headroom/small-cnn}.
Still to push to \texttt{origin/tmlr/submission}.

\end{document}
